\documentclass[10pt]{article}
\usepackage[letterpaper,margin=0.72in]{geometry}
\usepackage[T1]{fontenc}
\usepackage{lmodern,microtype,xcolor,enumitem,hyperref}
\definecolor{navy}{HTML}{08111F}
\definecolor{blue}{HTML}{2677FF}
\definecolor{muted}{HTML}{5D6B80}
\hypersetup{colorlinks=true,urlcolor=blue}
\setlist[itemize]{leftmargin=1.2em,itemsep=2pt,topsep=3pt}
\pagestyle{plain}
\begin{document}
{\color{blue}\small\bfseries SUMMER ENGINEERING PROJECT --- 2026}
\vspace{0.4em}
{\color{navy}\LARGE\bfseries Autonomous Lane Keeping on Tata Punch EV}\\
{\color{muted}\large Vehicle dynamics, perception, and closed-loop lateral control}

\section*{Project overview}
A summer 2026 engineering program centered on a lane-keeping architecture for a Tata Punch EV. The work links lane perception, path representation, lateral-error estimation, steering control, safety supervision, and repeatable validation.

\section*{Engineering problem}
A useful lane-keeping prototype must convert uncertain road observations into bounded steering commands while respecting vehicle dynamics and preserving a clear path to safe disengagement.

\section*{System architecture}
\begin{itemize}
\item Perception: identify lane boundaries or a drivable corridor from camera data.
\item State estimation: calculate lateral displacement, heading error, curvature, and confidence.
\item Control: compare geometric and model-based steering laws, with saturation and rate limits.
\item Safety: require manual override, confidence thresholds, speed limits, and fault logging.
\end{itemize}

\section*{Engineering focus}
\begin{itemize}
\item Bicycle-model reasoning for lateral motion and steering response.
\item Controller tuning across speed, curvature, sensor noise, and actuator delay.
\item Simulation-first validation before any controlled physical-vehicle experiment.
\end{itemize}

\section*{Verification plan}
\begin{itemize}
\item Track RMS and peak lateral error, heading error, steering rate, and disengagements.
\item Use scripted straight, constant-radius, S-curve, dropout, and delay scenarios.
\item Document assumptions and keep public-road deployment outside prototype scope.
\end{itemize}

\section*{Status and evidence boundary}
Active summer 2026 design brief; implementation details and measured results should be added only after logged tests.

\vfill
\small\color{muted} V. Narayan Kushwaha --- Summer Engineering Portfolio 2026
\end{document}
